\section{Lean formalization}\label{sec:formalization}

The Lean~4 formalization~\cite{lean4} follows the mathematical argument from the definition
of multiplication to the lower bound. It includes the finite
computations that establish the quotient premises, as well as the
geometric and algebraic deductions. In particular, the statements connecting quotient tensors to finite
integer systems are proved together with the certificate exclusions,
so the same formal argument covers their mathematical interpretation
and their application to the main theorem.

\subsection{Mathematical statements}
The field is represented by \code{ZMod 2}, and each factor is an
actual $3\times3$ matrix. The entrywise definition
of~\eqref{eq:tensor} is proved equivalent to the algorithm
identity~\eqref{eq:algorithm} for every pair of input matrices.
The final declaration \code{bilinear\_mul\_requires\_21} states that
every such algorithm has at least $21$ products. Zero summands are
allowed in the quantified family; no minimality, distinctness, or
symmetry condition is assumed.

The tensor-rank statement is \code{QiushiMatmul.rank\_ge\_21}.
Its proof combines the structural argument with a closed collection of
finite premises: bound $19$ for the three line types and the eight
all-high plane types, and bound $17$ for the fourteen normalized
affine-plane spans. The explicit upper-bound witness then gives
\code{rank\_between\_21\_and\_23}. These declarations formalize
Theorem~\ref{thm:main} and its stated interval.

\subsection{Finite proofs and reusable components}
The integer certificates of Section~\ref{sec:finite} are expanded
into Lean proofs. Each leaf proves a contradiction by an exact
nonnegative combination of inequalities. Each internal node applies an
exhaustive integer split. Quotient membership and the source rank
bounds are proved before these certificates are applied to a
hypothetical decomposition. Thus the final theorem derives its finite
premises within the same proof system as the saturation argument.

The development uses Mathlib's finite fields, matrices, subspaces,
linear maps, dimension theory, and rank
identities~\cite{mathlib}. The problem-specific components include
quotient/annihilator equivalence, occupation extraction, symmetry
transport, affine-coset geometry, and the multiplication-specific
product formula. Rank-additive saturation is formulated over any
field. It can be reused independently of the binary finite data.
The complete frozen-table theorem proves all $496$ representative
bounds and transports them to arbitrary subspaces; its logical place
is after the main theorem, as in
Proposition~\ref{prop:global}.

\subsection{Kernel verification}
The final theorems and their full transitive proof dependencies passed
Lean's unchanged kernel replay procedure from empty
environments~\cite{lean-replay}. Their axiom dependencies are confined to
\[
 \{\code{propext},\ \code{Classical.choice},\ \code{Quot.sound}\}.
\]
The main theorem uses no additional axiom for a quotient bound and no
external solver-success assumption. Its finite proofs are checked
proof terms, not trusted outputs of the certificate generators.

Appendix~\ref{app:formalization} lists the principal interfaces, the
pinned dependencies, and the replay groups.
The accompanying formalization package~\cite{qiushi-lean} contains the
sources, certificate data, generators, and commands for reproduction.
Neither a running Qiushi Engine nor access to a language model is
needed to check the proof.
